\documentclass[11pt]{article}
\usepackage{fontspec}
\usepackage[french,english]{babel}
\usepackage{amsmath}
\usepackage{amssymb}
\usepackage{amsthm}
\usepackage{mathrsfs}
\usepackage{tikz-cd}
\usetikzlibrary{babel}
\usepackage{enumitem}

\newcommand{\printedpage}[1]{\marginpar{\footnotesize p.~#1}}
\newcommand{\msp}[1]{\marginpar{\footnotesize [#1]}}
\newcommand{\M}{\mathcal{M}}
\newcommand{\Hom}{\mathrm{Hom}}
\newcommand{\HO}{\mathrm{HO}}
\newcommand{\Ob}{\mathrm{Ob}\,}
\newcommand{\Fl}{\mathrm{Fl}\,}
\newcommand{\Filt}{\mathrm{Filt}}
\newcommand{\card}{\mathrm{card}\,}
\newcommand{\rf}{\mathrm{f}}

\begin{document}

\printedpage{39}
% head: A. Grothendieck : Dérivateurs, Ch. XIII, Édition des Universités de Montpellier II et Paris VII
% foot: Version 25 janvier 2004

This therefore implies that the functor on $\HO_W(I)$ $F \longmapsto \Hom_{\HO_{W}(I')}(u^*F,G')$ is, if not representable, at least a term-by-term quotient of the functor represented by $u_*G'$ (\footnote[37]{But it is not clear, without an additional hypothesis on the situation, that this representable functor can be chosen in functorial dependence on $G'$, viewed in $\HO_W(I')$.}).

\msp{page 70}

But suppose now that $u$ is admissible. Then I am going to define the arrow in the original direction $(*)$ (p. 68) as follows: I choose an arrow
\[
i' \ : \ u^*F \longrightarrow F' \qquad \text{with } i' \in TC \subset C \cap W,\ F' \in \Ob \M(I')_{\rf}
\]
(so in fact $F' \in \Ob \M(I')_{\mathrm{cf}}$, since $u^*F$ is cofibrant), and by adjunction the arrow $i'$ corresponds to
\[
i \ : \ F \longrightarrow u_*F'.
\]
I consider the diagram (commutative!)
\[
\begin{tikzcd}[column sep=large, row sep=large]
\Hom(F',G') \arrow[r, "\alpha"] \arrow[d, "{\scriptstyle \lambda:f'\,\longmapsto\, f'\circ i'}"] & \Hom(u_*F',u_*G') \arrow[d, "{\scriptstyle \mu:f\,\longmapsto\, f\circ i}"] \\
\Hom(u^*F,G') \arrow[r, "\text{adj.}"] & \Hom(F,u_*G')\ .
\end{tikzcd}
\]

We note that we still have $F'$ cofibrant (since $u^*F$ is and $i' \in TC \subset C$), so $\Hom_W(F',G')$ is a quotient of $\Hom(F',G')$, and $\lambda$ induces an isomorphism
\[
\Hom_W(F',G') \xrightarrow{\ \sim\ } \Hom_W(u^*F',G') % [sic] u^*F' -> u^*F, the target of lambda is Hom(u^*F,G')
\]
(since $i' \in TC \subset W$). Since $F'$, $G' \in \Ob \M(I')_{\rf}$ and $u_*^{\rf}$ is compatible with $W(I')_{\rf}$, $W(I)_{\rf}$, it follows that the horizontal arrow $\alpha$
\msp{page 71}
is compatible with the equivalence relations defined by
\[
\Hom(F',G') \longrightarrow \Hom_{W(I')}(F',G') \quad \text{resp.} \ \ \Hom(u_*F',u_*G') \longrightarrow \Hom_{W(I)}(u_*F',u_*G')\ .
\]
And the same holds for the arrow $\mu : f \longmapsto f \circ i$, trivially. So the adjunction arrow $\Hom(u^*F,G') \longrightarrow \Hom(F,u_*G')$ is compatible with the relations $R$, $R'$, which finishes proving that $(**)$, page 68, induces a bijection $(*)$, page 68.

This therefore shows

\begin{enumerate}
\item[1°)] that for a fixed object $G'$ in $\HO_W(I')$ (which we can always suppose, up to isomorphism, to come from $\M(I')_{\rf}$), the functor
\[
F \ \longmapsto\ \Hom_{\HO_{W}(I')}(u^*_W F,G')
\]
is representable -- \textit{i.e.} the functor $u^*_W$ admits a left adjoint, and
\item[2°)] that this left adjoint is computed as $u_*(G')$ (when the argument $G'$ is taken fibrant) -- since the homomorphism $({*}{*}{*})$, page 69, is \textit{functorial in $F$}. (A posteriori it follows that it is also functorial in $G'$ \dots)
\end{enumerate}

\printedpage{40}
% head: A. Grothendieck : Dérivateurs, Ch. XIII, Édition des Universités de Montpellier II et Paris VII
% foot: Version 25 janvier 2004

\msp{page 72}

The last statement in 2° of the theorem-scholie follows from the computation of $u_*^W$, taking into account that $I'/s \longrightarrow I'$ is a local isomorphism, so the inverse image functor sends
\[
\M(I')_{\rf} \longrightarrow \M(I'/s)_{\rf},
\]
so if $G'$ is in $\M(I')_{\rf}$,
\[
u_*(G')(s) \ \simeq\ p_{s\,*}(G'|I'/s)
\]
is a computation of $p^W_{s\,*}(G'|I'/s)$.

This completes the proof of the theorem-scholie on direct images.

\textbf{Comments.}

1) This theorem seems to me far from satisfactory, it is a makeshift. I had hoped, with the hypotheses of paragraph 3, that we would have the existence of $u_*^W$ for \textit{every} arrow $u : I' \longrightarrow I$ in Cat. I now doubt that this is the case, but have not yet succeeded in making a counterexample. The most important test cases for the question of admissibility of a category $I$ seem to me the following:

\begin{enumerate}
\item[\textit{a})] $I$ an infinite discrete category (say countable). The question here comes down to this: let $f_i : X_i \longrightarrow Y_i$ be in $W$, with $X_i$, $Y_i$ fibrant, \textit{i.e.} in $F$ \texttt{[\textit{i.e.}} $X_i \longrightarrow e$\texttt{,} $Y_i \longrightarrow e$ \texttt{in} $F$\texttt{]}. Then is $f = \prod f_i : \prod X_i \longrightarrow \prod Y_i$ in $W$?

\msp{page 73}

\item[\textit{b})] I the category % [sic] I (upright roman) -> $I$, italic everywhere else for the category
\[
0 \longrightarrow 1 \longrightarrow 2 \longrightarrow 3 \longrightarrow \cdots,
\]
\textit{i.e.} the ordered category associated with the ordered set $\mathbf{N}$, for the usual order relation $m \leq n$. Thus $\M(I)$ consists of the \textit{inverse} systems in $\M$, indexed by $\mathbf{N}$. One easily proves (proceeding as in the corollary to the lemma p. 63) that $\M(I)_{\rf}$ consists of the inverse systems
\[
X_0 \longleftarrow X_1 \longleftarrow X_2 \longleftarrow \cdots \longleftarrow X_n \xleftarrow{\ p_n\ } X_{n+1} \longleftarrow \cdots
\]
such that the $X_i$ are in $\M_{\rf}$, and the $p_n$ in $F$. The question is therefore: if $X_\bullet$, $Y_\bullet$ are fibrant and $f_\bullet : X_\bullet \longrightarrow Y_\bullet$ is such that $f_i \in W$ for $i \in \mathbf{N}$, then is $\varprojlim f_i : X_\infty \longrightarrow Y_\infty$ in $W$?

If the answer is affirmative for $\mathbf{N}$ (for a $W$ and a given Quillenization), it % [sic] reponse -> réponse, accent missing
is also so for every countable discrete category, as one sees at once. But I do not see how to deduce even \textit{a}) from the axioms, say for a closed Quillen triple (closed model category) such that $\M$ is moreover stable under filtering $\varinjlim$, that these are exact, and that there is
\printedpage{41}
% head: A. Grothendieck : Dérivateurs, Ch. XIII, Édition des Universités de Montpellier II et Paris VII
% foot: Version 25 janvier 2004
\msp{page 74}
a small set of generators in $\M$ (so $\M$ is accessible), that $W$ is stable under filtering $\varinjlim$. I should ask \textsc{Thomason} whether he knows a counterexample. % [sic] connait -> connaît, circumflex missing

\item[\textit{c})] The two preceding examples were ordered categories. As a typical example of a non-ordered category, where I do not see what happens, I point out $I = B_G$, $G$ a group (so a single object, with an automorphism group $G$). Also the case of $\boldsymbol{\Delta}[1] = \{\boldsymbol{\Delta}_0, \boldsymbol{\Delta}_1\}$ (full subcategory of $\boldsymbol{\Delta}$) is not clear to me.
\end{enumerate}

2) One would like to be able to treat the existence of $u_!^W$ dually, for want of better. Under the starting hypothesis of paragraph 4 on $(\M,W)$, this therefore leads to a change of tactics. This time, we shall have to assume the existence of a closed Quillenization that is \textit{right exact} for $W$ (which is more demanding, it seems, than such a left exact structure, in the case under consideration), so
\[
TF \ = \ F \cap C,
\]
\msp{page 75}
and moreover also $C$, $TC$ such that we have
\[
C \ = \ \widetilde{C_0}, \qquad TC \ = \ \widetilde{TC_0}
\]
with $C_0$, $TC_0$ small (\footnote[38]{One then says that $C$, $TC$ (as right closed subsets \texttt{[rather on the left]}) are ``accessible''.}). We shall now set
\[
\begin{aligned}
F(I) \ &= \ \{f \in \Fl \M(I) \ \mid \ f_s \in F \quad\ \ \forall\, s \in \Ob I\}\\
TF(I) \ &= \ \{f \in \Fl \M(I) \ \mid \ f_s \in TF \ \ \forall\, s \in \Ob I\},
\end{aligned}
\]
then
\[
\begin{aligned}
C(I) \ &= \ TF(I)^*\\
TC(I) \ &= \ F(I)^*\ .
\end{aligned}
\]
The hope (which I shall have to check later) is that in this way we obtain a right-uniform Quillenization of $W$, over the whole of Cat. The only problem is the validity of the factorization axioms for $f \in \Fl \M(I)$ as $f = pi$, with $i \in TC$, $f \in F$, and also with $i \in C$, $f \in TF$. % [sic] f \in F, f \in TF -> p \in F, p \in TF, the factorisation is f = pi
I have good hope that this is true -- this reduces essentially to seeing that $(C(I),TF(I))$ and $(TC(I),F(I))$, say $(\Phi(I),\Psi(I))$, satisfy, in addition to $\Phi(I) = \Psi(I)^*$, $\Psi(I) = \Phi(I)_*$ (\textit{i.e.} $\Psi(I)$ is right closed), and that $\Phi(I)$ is accessible (just like $\Phi = \Phi(e)$). If this works, we shall thus have the existence
\printedpage{42}
% head: A. Grothendieck : Dérivateurs, Ch. XIII, Édition des Universités de Montpellier II et Paris VII
% foot: Version 25 janvier 2004
\msp{page 76}
of $u_!$ for $u$ a functor between finite ordered categories, and under slightly more general conditions. (E.g. $u : I' \longrightarrow I$ with $I$ ordered, $I'$ ordered and with underlying space noetherian and sober.) Perhaps this time, the hypothesis that $W$ is stable under filtering $\varinjlim$ will be of use after all! Thus, the case of discrete categories, and of left filtering categories such as $\mathbf{N}^\circ$, will no longer pose any problems because of this. So there is hope! Every ordered set which is a filtering $\varinjlim$ of \textit{closed} noetherian and sober ordered subsets, is then $W$-admissible on the right. But what about categories of type $\Gamma = B_G$ e.g. ? ? ?

\msp{page 77}

\section*{7\quad Cardinal filtrations and accessible subsets of a big category}

I refer to the theory of accessible functors and categories in SGA 4, I, §9 (p. 138-178). I give myself a big category (always a $\mathfrak{U}$-category) $E$ with a \textit{cardinal filtration} (increasing)
\[
\Bigl(\Filt^{\pi}(E)\Bigr)_{\pi \geq \pi_0}, \qquad \pi_0 \text{ a fixed infinite cardinal,}
\]
by the strictly full subcategories $\Filt^{\pi}(E)$. We have \texttt{[SGA 4, I, 9.12]}

\begin{enumerate}
\item[\textit{a})] $\Filt^{\pi}(E)$ essentially small, \textit{i.e.} equivalent to a small category (\textit{i.e.} to a category belonging to $\mathfrak{U}$).
\item[\textit{b})] $E$ satisfies $L_{\pi_0}$ (\textit{i.e.} $E$ is stable under $\varinjlim$ of type $I$, when $I$ is an ordered set large relative to $\pi_0$), and $\Filt^{\pi}(E)$ is stable under filtering direct limits indexed by $I$ large relative to $\pi_0$, and such that $\card I \leq \pi$.
\item[\textit{c})] For every $\pi \geq \pi_0$, and for every $X \in \Ob E$, we have
\[
X \ \simeq\ \varinjlim_{i \in I} X_i,
\]
where the $X_i$ \texttt{[are]} in $\Filt^{\pi}(E)$, and where $I$ is an ordered set large relative to $\pi$. Moreover, if $X \in \Ob \Filt^{\pi'}(E)$, with $\pi' > \pi$, we can take $I$ such that $\card I \leq \pi'^{\pi}$ (so if $\pi' = 2^{c'}$ with $c' \geq \pi$, whence $(\pi')^{\pi} = (2^{c'})^{\pi} = 2^{c'\pi} = 2^{c'} = \pi'$, we shall have $\card I \leq \pi'$). (\textbf{NB} If we assume the generalized continuum hypothesis, we shall always have $\pi'^{\pi} = \pi'$, hence $\card I \leq \pi'$.) (\footnote[39]{NB \textit{b}) and \textit{c}) imply that $E = \bigcup_{\pi \geq \pi_0} \Filt^{\pi}(E)$.})
\end{enumerate}

\printedpage{43}
% head: A. Grothendieck : Dérivateurs, Ch. XIII, Édition des Universités de Montpellier II et Paris VII
% foot: Version 25 janvier 2004

\msp{page 78 crossed out}

\msp{page 79}

Let us consider the representation in \textit{c}),
\[
X \ = \ \varinjlim_{I} X_i \qquad\qquad
\left\{
\begin{aligned}
& X_i \in \Filt^{\pi},\\
& I \text{ large relative to } \pi.
\end{aligned}
\right.
\]
Let
\[
\tilde{I} \ = \ \{I' \subset I \mid \card I' \leq \pi,\ \ I' \text{ large relative to } \pi_0\},
\]
endowed with the inclusion order relation. For
\[
I' \ \in\ \tilde{I}, \qquad \text{let}\ \ X_{I'} = \varinjlim_{I'} X_i\ ,
\]
this limit exists by $L_{\pi_0}$ (hypothesis \textit{b}), it is in $\Filt^{\pi}$ by \textit{b}), so we have a direct system
\[
(X_{I'})_{I' \in \tilde{I}}, \qquad X_{I'} \in \Filt^{\pi},
\]
with
\[
\card \tilde{I} \ \leq\ (\card I)^{\pi},
\]
and if $X \in \Filt^{\pi'}$ $(\pi' > \pi)$, taking
\[
\card I \ \leq\ \pi'^{\pi},
\]
we shall have $\card \tilde{I} \leq (\pi'^{\pi})^{\pi} = \pi'^{\pi^2} = \pi'^{\pi}$, hence
\[
\card \tilde{I} \ \leq\ \pi'^{\pi}.
\]
I say that $\tilde{I}$ is itself also \textit{large relative to $\pi$,} just like $I$. Indeed,
\msp{page 80}
let
\[
(I_\alpha)_{\alpha \in A}, \qquad I_\alpha \in \tilde{I}, \qquad \card A \ \leq\ \pi,
\]
let us prove that there exists $I' \in \tilde{I}$ which is an upper bound for the $I_\alpha$, \textit{i.e.} which contains $I'_0 = \bigcup_{\alpha \in A} I_\alpha$. We shall have $\card I'_0 \leq \pi^2 = \pi$ (since $\card A \leq \pi$ and $\card I_\alpha \leq \pi$ for every $\alpha$), and it suffices to see this:

\texttt{[In fact, the assertion is obvious: since} $\card I'_0 \leq \pi$\texttt{, and since} $I$ \texttt{is large relative to} $\pi$\texttt{,} $I'_0$ \texttt{admits an upper bound} $i'_0$ \texttt{in} $I$\texttt{, and one can take} $I' = I'_0 \cup \{i'_0\}$\texttt{.]}

\textbf{Lemma 1:} \textit{Let $\pi_0$, $\pi$ be two cardinals, with $\pi \geq \pi_0$, $I$ an ordered set large relative to $\pi$, and $I'_0$ a subset of $I$ of cardinality $\leq \pi$. Then there exists a subset $I'$ of $I$ satisfying the conditions:} (\footnote[40]{Provided one assumes $\pi > \pi_0$, and moreover $\pi^{\pi_0} = \pi$ (e.g. $\pi = 2^c$ with $c \geq \pi_0$).})

\begin{enumerate}
\item[a)] $I' \supset I'_0$.
\item[b)] $\card I' \leq \pi$.
\end{enumerate}

\end{document}
